%-------------------------
% Cover Letter -- Luma AI, Agent Behavior Designer
% Min-Han Li
%------------------------
\documentclass[letterpaper,11pt]{article}

\usepackage[margin=1in]{geometry}
\usepackage[hidelinks]{hyperref}
\usepackage[T1]{fontenc}

\pagestyle{empty}
\setlength{\parindent}{0pt}
\setlength{\parskip}{10pt}

\begin{document}

Dear Luma Team,

I'm applying for the Agent Behavior Designer role. Treating an agent's behavior, not just its code, as the thing you design is exactly how I already work, and ``prose functions as code'' is how I think.

My path runs from chemical engineering to CS to getting language models to actually behave. At ByteDance I built an on-call diagnostic agent for 200+ engineers: I designed its prompt stacks, gave it tools to call, and spent most of my time reading traces to find why it misbehaved and rewriting prompts until it didn't. That loop is the part I love.

On Pebble, a friend's CBT mental-health startup, I handle the agent's software. I built it on a skill-as-node architecture with per-node evaluators and a crisis path that cannot be bypassed, because when an agent talks to a vulnerable person, ``usually fine'' is not enough. My cognitive science and HCI background keeps me asking what the user actually needs.

I have a full-time offer to return to ByteDance as a software engineer, but this role is the one I actually want. It matches where I'm taking my career: shaping how agents behave, not just building backends. Luma puts agents into the real world inside a tool people use, and I'd love to bring that obsession to your canvas.

Thank you for your time and consideration.

\vspace{6pt}

Sincerely, \\
Min-Han Li \\
minhan.li.wesley@gmail.com $|$ +1-781-475-7402

\end{document}
